\documentclass[../main.tex]{subfiles}

\begin{document}

\section{Chapter 8: Interaction Techniques}
\label{sec:chapter8-interaction-techniques}

\subsection{Techniques / Principles Applied}

\begin{itemize}[leftmargin=*]
\item Overview first, then filter, then details-on-demand.
\item Coordinated views that update several visuals from the same selection.
\item Local page-level filters instead of one global interaction state.
\item Power BI slicers and linked visuals as the main dashboard interaction
tools.
\end{itemize}

\subsection{How Applied in the Dashboard}

The dashboard was built in Power BI and uses interaction mainly to extend a
guided policy narrative into subgroup exploration. Users first see overview
states on pages such as Account Ownership \& Access Landscape and Financial
Health \& Resilience, then narrow the analysis through slicers, and finally use
details-on-demand in the live dashboard to inspect exact values.

\dashboardfigure{0.92}{fig-page3-filter-bar.png}
{List of filters on the Account Ownership \& Access Landscape page.}
{fig:interaction-page3-filter-bar}

Figure~\ref{fig:interaction-page3-filter-bar} shows the main interaction design
on Page 3: the page starts with a national overview, then lets users filter by
urbanicity, income group, gender, education level, and workforce status.
Within the live Power BI dashboard, those slicers coordinate the KPI cards, the
Sankey diagram, and the exclusion-barrier chart so that all three reflect the
same selected cohort.

\dashboardfigure{0.82}{fig-page3-account-landscape-overview.png}
{Default Account Ownership \& Access Landscape state visible in the static PDF export.}
{fig:page3-interaction-states}

The exported PDF confirms the placement of the filters and the default overview
state. It also supports the claim that the dashboard uses local page-level
interaction rather than one global control state shared across every page.

\subsection{Notes / Adjustments}

Because the available source is a static PDF export, this report cannot verify
live hover tooltips, zoom behavior, animation, or the exact values produced by
every filtered state. The chapter therefore describes the dashboard as designed
for filtering, coordinated updates, and details-on-demand without claiming
interactive evidence that is not directly visible in the exported file.

\end{document}
